\begin{textsample}{POS Dim 4 – 2018\_c6 – Score 15.00 – 2018\_c6/enade\_2018\_admin\_6a5e0b186f.txt}  \label{ex:f4_pos_003}
CAPÍTULO 6 ANÁLISE TÉCNICA DA PROVA Este capítulo tem por objetivo apresentar o desempenho dos estudantes concluintes de Administração no Enade/2018.
Para isso, foram calculadas as estatísticas básicas da prova em seu todo ( seção 6.1.1 ), bem como as estatísticas relacionadas aos Componentes de Formação Geral ( seção 6.1.2 ) e de Conhecimento Específico da Área ( seção 6.1.3 ).
Dadas as suas características, foram analisadas, em separado, as questões objetivas ( seção 6.2 ) e as questões discursivas ( seção 6.3 ).
Para as questões objetivas, foram disponibilizados os índices de facilidade e de discriminação Ponto-Bisserial, também em separado, para os Componentes de Formação Geral ( seção 6.2.1 ) e de Conhecimento Específico ( seção 6.2.2 ).
De cada componente, uma das questões foi escolhida para exemplificar a análise gráfica, relacionando as alternativas escolhidas pelos estudantes ( inclusive o gabarito ) com o número de acertos no componente.
No Anexo I, apresenta -se a íntegra da análise gráfica para todas as questões objetivas.
Para cada uma das questões discursivas, os conteúdos dos tipos mais comuns de respostas dos estudantes são apresentados e comparados com o padrão de resposta esperado ( ver Anexo VIII com o padrão de respostas ).
Tomando -se como base as duas questões discursivas do Componente de Formação Geral, na seção 6.3.1.6, são apresentados comentários sobre a correção das respostas em relação à Língua Portuguesa.
Nas tabelas, constam as seguintes estatísticas das notas20 : média do desempenho na prova, erro padrão da média, desvio padrão, nota mínima, mediana e nota máxima para cada um de seus componentes.
Tais estatísticas contemplam o total de estudantes concluintes da Área de Administração inscritos e presentes à prova do Enade/2018, tendo em vista agregações, ou por Grandes Regiões e o país como um todo, ou por Categoria Administrativa, Organização Acadêmica da IES e Modalidade de Ensino.
Em relação aos gráficos de distribuição de notas, o intervalo considerado foi de 10 unidades, aberto à esquerda e fechado à direita, com exceção do primeiro intervalo, [ 0 ; 10 ], fechado em ambos os extremos.
Para os gráficos de distribuição das notas das questões discursivas, foram consideradas mais duas categorias : questão em branco21 e nota zero. 20 Uma definição dessas estatísticas pode ser encontrada no Glossário. 21 Nesse grupo estão incluídas também as respostas classificadas como nulas ou desconsideradas. 6.1 ESTATÍSTICAS BÁSICAS DA PROVA Esta seção apresenta estatísticas selecionadas e histogramas da nota geral ( 6.1.1 ) e de cada componente : Formação Geral ( 6.1.2 ) e Conhecimento Específico ( 6.1.3 ).
São, também apresentadas estatísticas selecionadas de subpopulações, caracterizadas por Grande Região, Categoria Administrativa, Organização Acadêmica e Modalidade de Ensino. 6.1.1 Estatísticas Básicas Gerais Na Tabela 6.1, são apresentadas as Estatísticas Básicas da prova, por Grande Região, dos estudantes concluintes de Administração.
A Média das notas da prova, como um todo, ( nas seções seguintes serão analisados os Componentes de Formação Geral e de Conhecimento Específico ) foi 38,5, sendo que os alunos da região Centro-Oeste obtiveram a Média mais baixa ( 37,3 ), e os da região Sudeste obtiveram a Média mais alta ( 39,3 ).
As demais médias foram : 37,9, na região Norte, 38,9, na região Nordeste e 37,5, na região Sul.
O Desvio padrão para o Brasil, como um todo, foi 13,2, sendo o maior Desvio padrão encontrado na região Sudeste ( 13,3 ), e o menor, na região Norte ( 12,3 ), indicando uma dispersão um pouco menor das notas dessa última região.
A região que obteve a maior nota Máxima foi a Nordeste ( 92,8 ), ao passo que a região que atingiu a menor nota Máxima foi a Norte ( 76,8 ).
A Mediana do Brasil, como um todo, foi 37,8, sendo a maior Mediana obtida na região Sudeste ( 38,7 ), e a menor obtida na região Centro-Oeste ( 36,5 ).
A nota Mínima do Brasil foi obtida nas regiões Nordeste, Sudeste, Sul e na Centro-Oeste, onde esta foi zero.
A maior nota Mínima, entre as cinco Grandes Regiões, foi 4,8, na região Norte, as demais regiões obtiveram nota Mínima igual à zero.
Considerando -se as notas segundo Grande Região, observa -se que não existe diferença estatisticamente significativa ao nível de 95 \% entre as Médias, obtidas nas regiões Norte, Sul e Centro-Oeste.
O comportamento das notas dos estudantes de todo o Brasil pode ser observado no Gráfico 6.1, que apresenta um histograma com a distribuição das mesmas : uma distribuição unimodal com a moda no intervalo ( 30 ; 40 ].
Na Tabela 6.2, são apresentadas informações referentes à Média da nota final, desagregadas por Categoria Administrativa, por Organização Acadêmica e por Modalidade de Ensino.
A Média das notas da prova, como um todo, foi 38,5.
Em relação à Categoria Administrativa, os estudantes das IES Públicas obtiveram Média mais alta ( 44,7 ) que a Média nacional, e os das IES Privadas obtiveram Média mais baixa ( 37,5 ).
Observa -se que existe diferença estatisticamente significativa entre as médias das notas das IES Públicas e as das Privadas.
A diferença entre as médias das regiões Sudeste e Centro-Oeste ( 2,0 ), a maior e a menor Média, é inferior à diferença entre a das IES Públicas e Privadas ( 7,3 ), caracterizando -se uma maior diversidade administrativa do que regional.
O Desvio padrão para as IES Públicas ( 14,0 ) foi superior ao do Brasil, como um todo ( 13,2 ), indicando uma dispersão um pouco maior das notas nesta Categoria Administrativa.
No tocante à Organização Acadêmica, as Universidades e os CEFET/IFET obtiveram Média mais alta que a nacional ( 39,4 e 44,6, respectivamente ).
A Média dos Centros Universitários e a das Faculdades foram menores do que a nacional ( 38,1 e 37,1, respectivamente ).
Constata -se que existe diferença estatisticamente significativa, ao nível de 95 \%, entre a Média mais alta, 44,6 dos CEFET/IFET, e a dos demais tipos de Organização Acadêmica.
A Média da Modalidade Educação a Distância ( 35,4 ) foi inferior à Média da Educação Presencial ( 39,5 ), e a diferença entre elas é estatisticamente significativa, ao nível de 95 \%. 6.1.2 Estatísticas Básicas no Componente de Formação Geral Na Tabela 6.3, são apresentadas as Estatísticas Básicas em relação ao componente da prova que avalia a Formação Geral dos estudantes concluintes.
Os alunos de todo o Brasil obtiveram desempenho médio de 43,6.
Quanto à variabilidade, o Desvio padrão das notas dos estudantes do Brasil, como um todo, foi 16,6.
A maior Média foi obtida nas regiões Nordeste e Sudeste ( 44,2 ), e a menor, na região Centro-Oeste ( 41,8 ).
As demais médias foram : 43,8, na região Norte e 42,7, na região Sul.
Já o maior Desvio padrão foi obtido nas regiões Nordeste, Sudeste e Centro-Oeste ( 16,7 ) e o menor, na região Norte ( 15,9 ).
Na região Sul o Desvio padrão foi 16,2.
A maior nota no Componente de Formação Geral da prova do Enade/2018 foi 96,4, obtida por, pelo menos, um aluno na região Sudeste.
A menor nota Máxima foi obtida na região Norte ( 91,4 ).
A Mediana do Brasil, como um todo, foi 43,0, sendo a menor Mediana encontrada na região Centro-Oeste ( 41,2 ), e a maior, na região Nordeste ( 43,8 ).
A nota Mínima nessa parte foi zero em todas as regiões.
Considerando -se as notas, segundo Grande Região, observa -se que não existe diferença estatisticamente significativa entre as Médias das notas do Componente de Formação Geral obtidas nas regiões Norte, Nordeste e Sudeste.
As diferenças entre as demais Médias são estatisticamente significativas.
No Gráfico 6.2, é possível verificar a avaliação do desempenho dos estudantes do Componente de Formação Geral, a partir do histograma da distribuição das notas correspondentes.
A distribuição é unimodal, com moda em ( 40 ; 50 ], intervalo imediatamente seguinte ao modal da distribuição de notas da prova, como um todo ( Gráfico 6.1 ).
Nota -se, ainda, que, no Gráfico 6.2, as notas apresentam maior dispersão do que as no Gráfico 6.1 ( distribuição das notas da prova ), confirmada pela comparação dos desvios padrões : 13,2 para a nota da prova, como um todo, e 16,6 para o Componente de Formação Geral.
Na Tabela 6.4, são apresentadas as informações referentes ao desempenho dos concluintes do Componente de Formação Geral, em diferentes agregações : Categoria Administrativa, Organização Acadêmica e Modalidade de Ensino.
Observa -se que existe diferença estatisticamente significativa entre as médias dos tipos de Categoria Administrativa : a maior Média obtida por estudantes de IES Públicas ( 50,9 ) e a menor, pelos de IES Privadas ( 42,3 ).
Nota -se que, também, há diferença estatisticamente significativa entre a maior das médias dos quatro tipos de Organização Acadêmica e as demais, o valor maior para os CEFET/IFET ( 49,8 ) e as demais : 45,1 para as Universidades, 42,8 para os Centros Universitários e 41,5 para as Faculdades, a menor Média.
No que se refere à Modalidade de Ensino, a Educação Presencial apresentou maior Média do que a Educação a Distância ( médias 44,4 e 41,2, respectivamente ).
Tais médias apresentam diferença estatisticamente significativa. 6.1.3 - Estatísticas Básicas do Componente de Conhecimento Específico Na Tabela 6.5, são apresentadas as Estatísticas Básicas referentes ao Componente de Conhecimento Específico da área de Administração.
A Média do desempenho dos alunos do Brasil, como um todo, foi 36,8.
A maior Média foi obtida na região Sudeste ( 37,6 ), e a menor, nas regiões Sul e Centro-Oeste ( 35,7 ).
As demais médias foram : 36,0, na região Norte e 37,1, na região Nordeste.
Quanto à variabilidade das notas, o Desvio padrão do Brasil, como um todo, foi 14,2, sendo o maior Desvio padrão observado na região Sudeste ( 14,3 ) e o menor, na região Norte ( 13,4 ).
Os demais desvios foram : 14,1, na região Nordeste, 14,2, na região Sul e 13,9, na região Centro-Oeste.
A Mediana das notas dos estudantes de todo o Brasil foi 36,6.
A maior Mediana ocorreu na região Sudeste ( 37,6 ), e a menor, na região Sul ( 35,1 ).
As demais medianas foram : 35,6 na região Norte, 36,8 na região Nordeste e 35,6 na região Centro-Oeste.
A nota Máxima do Brasil, como um todo, foi 95,2, sendo obtida por, pelo menos, um aluno na região Nordeste.
As demais notas máximas foram : 80,8, na região Norte, 92,0, na região Sudeste, 87,7, na região Sul e 84,1 na região Centro-Oeste.
A nota Mínima foi zero em todas as regiões.
Observa -se que não existe diferença estatisticamente significativa entre as três menores Médias, das notas do Componente de Conhecimento Específico, das regiões Norte, Sul e Centro-Oeste.
No entanto, existe diferença estatisticamente significativa entre as demais regiões.
Assim como nos Gráficos 6.1 e 6.2, no Gráfico 6.3, apresentado a seguir, é possível ser feita uma avaliação do desempenho de concluintes de Administração, em relação ao Componente de Conhecimento Específico, com um histograma da distribuição das notas correspondentes.
Esta também é uma distribuição unimodal, e o grupo modal é o ( 30 ; 40 ], o mesmo do grupo modal da prova, como um todo, e abaixo do grupo modal para a Formação Geral.
Na Tabela 6.6, é apresentada uma comparação dos resultados em relação à Categoria Administrativa, à Organização Acadêmica e a Modalidade de Ensino, agora levando -se em conta o desempenho de alunos do Componente de Conhecimento Específico da prova.
No que se refere à Organização Acadêmica, a maior Média foi a dos CEFET/IFET ( 42,9 ), vindo a seguir a das Universidades ( 37,5 ), a dos Centros Universitários ( 36,6 ) e, por fim, a das Faculdades ( 35,7 ).
O maior Desvio padrão, e acima do valor para o Brasil, como um todo, foi o dos CEFET/IFET ( 15,0 ).
As Universidades obtiveram a maior nota Máxima ( 95,2 ).
Os Centros Universitários obtiveram nota Máxima 85,8 ; as Faculdades, 91,2 e os CEFET/IFET obtiveram nota Máxima 85,2.
As medianas foram : 37,3 nas Universidades, 36,1 nos Centros Universitários e 35,3 nas Faculdades, a menor delas ; e 42,8, nos CEFET/IFET a maior.
A nota Mínima foi zero para todas as Organizações Acadêmicas, exceto para os CEFET/IFET ( 4,8 ).
Observa -se que existe diferença estatisticamente significativa, ao nível de 95 \%, no Componente de Conhecimento Específico entre as notas nos quatro tipos de Organização Acadêmica.
Quanto à Categoria Administrativa, observa -se um comportamento semelhante àquele da parte de Formação Geral e à prova, como um todo, ou seja, existe diferença estatisticamente significativa entre as médias das IES Públicas ( 42,6 ) e as das IES Privadas ( 35,8 ).
Neste caso, também, a maior Média foi obtida por alunos de IES Públicas de ensino.
Quanto ao comportamento das médias, no que toca à Modalidade de Ensino, há diferença estatisticamente significativa entre a Média da Educação Presencial ( 37,8 ) e a da Educação a Distância ( 33,5 ). 6.2 ANÁLISE DAS QUESTÕES OBJETIVAS Esta seção apresenta estatísticas selecionadas e histogramas das Questões Objetivas de Formação Geral ( 6.2.1 ) e de Conhecimento Específico ( 6.2.2 ).
São, também, apresentadas e comparadas as médias das subpopulações caracterizadas por Grande Região. 6.2.1 Componente de Formação Geral Na Tabela 6.7, são apresentadas as Estatísticas Básicas relativas às oito questões objetivas do componente da prova que abrange a Formação Geral dos estudantes.
A Média do Brasil foi 48,0.
A menor Média foi encontrada na região Centro-Oeste ( 46,5 ), e a maior, na região Nordeste ( 48,8 ).
As demais médias foram : 47,2, na região Norte, 48,6, na região Sudeste e 46,9, na região Sul.
O Desvio padrão do Brasil foi 22,6, sendo o maior Desvio padrão encontrado nas regiões Nordeste e Sudeste ( 22,7 ), e o menor, na região Norte ( 21,7 ).
Os demais desvios foram : 22,3, na região Sul e 22,6, na região Centro-Oeste.
As medianas do Brasil, como um todo, e das cinco regiões foram iguais a 50,0.
A nota Máxima 100,0 foi alcançada nas cinco Grandes Regiões.
As notas mínimas ( 0,0 ) foram iguais para todas as regiões.
Na Tabela 6.8, são apresentados o Índice de Facilidade e o Índice de Discriminação ( Ponto-Bisserial ) para cada uma das questões objetivas do Componente de Formação Geral.
Quanto ao Índice de Facilidade, foram usadas as seguintes cores para diferenciar o nível de dificuldade da questão :  Azul para as questões classificadas com índice Muito fácil ( \&gt ; =0,86 ), verde para as questões classificadas com índice Fácil ( 0,61 a 0,85 ), amarelo para as questões classificadas com Médio ( 0,41 a 0,60 ), vermelho para as questões classificadas com Difícil ( 0,16 a 0,40 ) e roxo para as questões classificadas com Muito difícil ( \&lt ; =0,15 ).
Já quanto ao Índice de Discriminação, foram usadas as seguintes cores para qualificar a questão :  As questões classificadas com índice Fraco receberam a cor vermelha ( \&lt ; =0,19 ), as classificadas com Médio receberam a cor amarela ( 0,20 a 0,29 ), as classificadas com Bom receberam a cor verde ( 0,30 a 0,39 ) e as classificadas com Muito bom ( \&gt ; =0,40 ) receberam a cor azul.
As questões objetivas do Componente de Formação Geral, segundo o Índice de Facilidade, foram assim avaliadas : das oito questões, nenhuma teve o Índice de Facilidade classificado como Muito fácil, e uma questão foi tida como Fácil, por ter índice de acertos 0,65.
Cinco questões foram consideradas com índice de dificuldade Médio, situando -se no intervalo entre 0,45 e 0,60 do Índice de Facilidade, ou seja, houve entre 45,0 \% e 60,0 \% de acertos, enquanto duas questões foram classificadas na categoria Difícil, situando -se no intervalo entre 0,27 e 0,40.
Nenhuma questão apresentou menos de 15 \% de acertos, razão pela qual não houve questão classificada como Muito difícil.
O Índice de Facilidade variou de 0,27 a 0,65 e o de Discriminação, de 0,42 a 0,50.
Todas as questões, quanto ao Índice de Discriminação, foram tidas como Muito bom.
Na Tabela 6.9, é apresentada a distribuição das questões, levando -se em conta, simultaneamente, a classificação dos dois índices.
As oitos questões que tiveram Índice de Discriminação Muito bom figuraram entre três níveis de dificuldade Fácil, Médio e Difícil : uma classificada na categoria Fácil ( a questão 2 ) do Índice de Facilidade, cinco, na categoria Médio ( questões 1, 3, 6, 7 e 8 ), e duas, na categoria Difícil ( questões 4 e 5 ).
Em particular, a questão 1 foi a que apresentou o maior poder discriminatório, com índice 0,50, e foi considerada Médio em termos de facilidade, com uma proporção de 0,49 acertos.
O máximo de acertos foi alcançado pela questão 2, com um Índice de Facilidade de 0,65.
As questões 4 e 5 foram consideradas Difícil, com índices de Facilidade 0,27 e 0,40, respectivamente.
Os índices de Discriminação foram Muito bom para as duas questões ( 0,49 e 0,48 respectivamente ).
No Gráfico 6.4, para exemplificar, analisa -se o comportamento da questão de número 4 de Formação Geral.
Trata -se de uma questão considerada Difícil, em relação à facilidade e a que obteve o segundo maior Índice de Discriminação dessa parte da prova.
Neste gráfico, cada uma das cinco curvas representa o percentual de respostas em determinada alternativa da questão, em função do número de acertos dos estudantes nessa parte da prova ( Formação Geral/Múltipla Escolha ), antes de possíveis eliminações pelo critério Ponto-Bisserial.
A curva em verde corresponde à alternativa B, a alternativa correta Na Tabela 6.9, é apresentada a distribuição das questões, levando -se em conta, simultaneamente, a classificação dos dois índices.
As oitos questões que tiveram Índice de Discriminação Muito bom figuraram entre três níveis de dificuldade Fácil, Médio e Difícil : uma classificada na categoria Fácil ( a questão 2 ) do Índice de Facilidade, cinco, na categoria Médio ( questões 1, 3, 6, 7 e 8 ), e duas, na categoria Difícil ( questões 4 e 5 ).
Em particular, a questão 1 foi a que apresentou o maior poder discriminatório, com índice 0,50, e foi considerada Médio em termos de facilidade, com uma proporção de 0,49 acertos.
O máximo de acertos foi alcançado pela questão 2, com um Índice de Facilidade de 0,65.
As questões 4 e 5 foram consideradas Difícil, com índices de Facilidade 0,27 e 0,40, respectivamente.
Os índices de Discriminação foram Muito bom para as duas questões ( 0,49 e 0,48 respectivamente ).
No Gráfico 6.4, para exemplificar, analisa -se o comportamento da questão de número 4 de Formação Geral.
Trata -se de uma questão considerada Difícil, em relação à facilidade e a que obteve o segundo maior Índice de Discriminação dessa parte da prova.
Neste gráfico, cada uma das cinco curvas representa o percentual de respostas em determinada alternativa da questão, em função do número de acertos dos estudantes nessa parte da prova ( Formação Geral/Múltipla Escolha ), antes de possíveis eliminações pelo critério Ponto-Bisserial.
A curva em verde corresponde à alternativa B, a alternativa correta para essa questão.
Observa -se que entre os estudantes com menor número de acertos, nessa parte do exame, a situação mais frequente foi a escolha da alternativa C ( em preto ), incorreta.
Por exemplo, entre os estudantes que acertaram três questões, 19,6 \% escolheram a alternativa E ( em vermelho ), 15,4 \% escolheram a alternativa D ( em laranja ), 12,5 \% escolheram a alternativa B ( em verde, gabarito ), 10,9 \% escolheram a alternativa A ( em azul ), e 41,5 \%, a C ( em preto ).
Dentre os que acertaram três respostas entre as questões de múltipla escolha de Formação Geral, 0,1 \% deixou a questão em branco e 0,1 \% marcou mais de uma alternativa, invalidando a questão.
À medida que o número de acertos aumenta, indicando desempenho melhor nessa parte da prova, aumenta, concomitantemente, a proporção de estudantes que selecionaram a alternativa correta B, atingindo 100 \% para os estudantes com oito acertos.
Essa análise permite verificar como a questão discriminou os grupos de desempenho, justificando -se o alto índice ( 0,49 ) obtido na questão.
Cumpre notar que não é possível inferir deste gráfico nem o índice de facilidade ( que seria uma média da proporção ponderada pela quantidade de alunos com cada uma das notas ), nem o índice de discriminação Ponto-Bisserial, por razão equivalente.
No caso extremo, no qual a grande concentração dos acertos dos alunos fosse abaixo de quatro, o índice de facilidade seria obrigatoriamente abaixo de 20,0 \% ( neste exemplo ).
Caso a concentração fosse em seis acertos ou mais, o índice seria obrigatoriamente acima de 50,0 \%.
Os gráficos relativos às demais questões de Formação Geral constam do Anexo I. 6.2.2 Componente de Conhecimento Específico Na Tabela 6.10, são apresentadas as Estatísticas Básicas em relação às questões objetivas do Componente de Conhecimento Específico da prova de Administração por Grande Região.
Nessa parte da prova, uma questão foi anulada pela Comissão Assessora da Área, a questão 23.
Além disso, como será discutido ainda nesta seção, oito questões objetivas do Componente de Conhecimentos Específicos não foram usadas no cômputo das notas por terem sido descartadas pelo critério do Ponto-Bisserial.
Assim, as notas foram calculadas com base em 18 das 26 questões objetivas válidas de Conhecimentos Específicos.
A Média do Brasil deste componente foi 40,0.
A menor Média foi observada na região Sul ( 38,8 ), e a maior, na região Sudeste ( 40,7 ).
O Desvio padrão de todo o Brasil foi 15,6, sendo o menor Desvio padrão encontrado na região Norte ( 14,9 ), e o maior, na região Sudeste ( 15,7 ).
A Mediana de todo o Brasil foi 38,9, o mesmo valor da Mediana encontrada em todas as regiões.
A nota Máxima da prova foi 100,00, obtida nas questões objetivas do Componente de Conhecimento Específico, por, pelo menos, um aluno das regiões Nordeste e Sudeste.
A nota Mínima foi zero em todas as regiões.
Na Tabela 6.11, são apresentados os Índices de Facilidade e Discriminação ( Ponto-Bisserial ) das questões objetivas do Componente de Conhecimento Específico, para os estudantes de Administração.
Para facilitar a diferenciação das questões, serão usadas as mesmas cores da Tabela 6.8 para as diferentes classificações dos Índices de Facilidade e de Discriminação.
Observa -se que uma questão foi anulada pela CAA e que as análises dizem respeito às 26 questões válidas.
A partir do Índice de Facilidade obtido, pode -se concluir que mais da metade das questões objetivas da prova foram considerados, pelo menos, Difícil : das 26 questões válidas, 17 foram classificadas como Difícil ou como Muito difícil.
Uma questão foi classificada como Fácil, e outras oito consideradas como Médio.
Nenhuma questão foi classificada como Muito fácil Já quanto aos índices de discriminação das questões objetivas do Componente de Conhecimento Específico da prova, obtém -se como resultado a seguinte classificação : nove das 26 questões foram consideradas boas, enquanto uma delas teve Índice de Discriminação Muito bom.
Assim, para 10 em 26 questões, os Índices de Discriminação foram Bom ou Muito bom.
Dentre as demais, oito delas foram classificadas como Médio, e outras oito, como Fraco, sendo 16, por conseguinte, a quantidade de questões nos dois patamares mais baixos de discriminação.
Constata -se, assim, que a prova – no que se refere ao Componente de Conhecimento Específico – possuía média capacidade de discriminar entre aqueles que dominam ou não o conteúdo.
O Índice de Facilidade variou de 0,10 a 0,61 e o de Discriminação, de 0,01 a 0,42.
Na Tabela 6.12, é apresentada a distribuição das questões levando -se em conta, simultaneamente, a classificação dos dois índices.
Dentre as questões que alcançaram os maiores índices de discriminação, uma delas foi classificada com Índice de Discriminação Muito bom, a de número 24.
Seu índice foi 0,42, e, quanto ao índice de Facilidade, foi classificada como Médio, com 46 \% dos estudantes marcando a opção correta.
A questão de número 10 foi a mais difícil entre as 26 questões específicas válidas, com baixo Índice de Facilidade, apenas 10 \% de acertos.
Essa questão apresentou poder discriminatório muito baixo, 0,13, o que comprova ter sido a mais difícil para os estudantes.
Destaca -se, também, a questão 9, com Índice de Facilidade 0,11, o que, em termos percentuais, corresponde a 11 \% de estudantes que responderam acertadamente.
Já 0,01 foi o seu Índice de Discriminação.
Tais questões foram, portanto, pelo critério Ponto-Bisserial, consideradas inadequadas.
Por isso, as questões 9 e 10, além das questões 11, 12, 29, 31, 34 e 35, foram eliminadas do cômputo da nota final.
A título de exemplo das análises do comportamento das questões objetivas, no Gráfico 6.5, analisa -se a questão 24 do Componente de Conhecimento Específico.
Essa questão foi considerada, pelas respostas dos estudantes avaliados na prova, como uma questão com nível de facilidade Médio, 0,46, ou seja, 46 \% dos estudantes assinalaram, acertadamente, a opção A, correspondente ao gabarito.
Como já comentado, seu Índice de Discriminação foi igual a 0,42, classificado como Muito bom, o maior valor de discriminação.
Neste gráfico, cada uma das cinco curvas representa o percentual de respostas em determinada alternativa da questão 24, em função do número de acertos dos estudantes nessa parte da prova, antes de possíveis eliminações de questões pelo critério Ponto-Bisserial.
A alternativa correta A, representada no gráfico pela curva em azul, foi escolhida em maiores proporções pelos alunos com desempenho melhor nessa parte da prova.
Já as alternativas incorretas, também denominadas distratores, foram selecionadas, principalmente, por aqueles com notas mais baixas.
Observa -se que a soma não é 100 \%, por causa das questões não respondidas ou com mais de uma opção marcada.
Aqueles com nota zero, na sua quase totalidade, deixaram esta questão em branco ou marcaram mais de uma alternativa, comportamento considerado inválido.
A proporção de alunos que selecionaram a resposta correta A aumenta gradativamente, chegando a atingir 100 \% para 21 acertos ou mais, enquanto a proporção dos que escolheram alternativas incorretas decai, a partir de três ou quatro acertos, em função do número de acertos nessa parte da prova.
Os gráficos relativos às demais questões do Conhecimento Específico constam do Anexo I. 6.3 ANÁLISE DAS QUESTÕES DISCURSIVAS Esta seção apresenta estatísticas selecionadas e histogramas das Questões Discursivas de Formação Geral ( 6.3.1 ) e Conhecimento Específico ( 6.3.2 ).
São, também, apresentadas e comparadas as médias de subpopulações, caracterizadas por Grande Região. 6.3.1 Componente de Formação Geral As análises dos resultados de desempenho dos estudantes de Administração, nas duas questões discursivas relativas à Formação Geral, encontram -se na Tabela 6.13 e no Gráfico 6.6.
Na Tabela 6.13, observa -se que a nota Média nesse conjunto de questões foi inferior à obtida nas objetivas.
Os estudantes de todo Brasil obtiveram, em Formação Geral, Média 48,0, nas questões objetivas e 37,1, nas questões discursivas.
Pode -se notar, também, que o Desvio padrão nesse conjunto de questões foi menor do que o obtido nas objetivas : 22,6, nas questões objetivas e 17,6, nas questões discursivas.
A maior Média foi obtida na região Norte ( 38,8 ), e a menor, na região Centro-Oeste ( 34,7 ).
A Mediana de todo o Brasil, neste componente, foi 39,0.
Nas regiões Sul e Centro-Oeste, a Mediana foi menor ( 38,0 e 37,0, respectivamente ), e, nas demais regiões foi maior que o valor para o Brasil : Norte ( 41,0 ), Nordeste e Sudeste ( 39,5 ).
A nota Máxima ( 93,0 ) foi obtida na região Sudeste, sendo 84,0, a Máxima na região Norte ; 92,5, na Nordeste ; 89,5, na Sul ; e 88,0, na região Centro-Oeste.
A nota Mínima ( 0,0 ) foi a mesma em todas as regiões do Brasil.
No Gráfico 6.6, está representada a distribuição das notas nas questões discursivas do Componente de Formação Geral.
A moda dessa distribuição ocorre no intervalo ( 40 ; 50 ], seguida pelo intervalo ( 30 ; 40 ].
Destaca -se, também, o intervalo [ 0 ; 10 ], com distribuição de 10,1 \% do total de notas, sendo que no intervalo [ 0 ; 10 ], inclui -se, além da nota zero, a frequência de alunos que deixaram esse tipo de questão em branco.
Na sequência, os resultados verificados para cada uma das questões discursivas de Formação Geral serão apresentados, estabelecendo -se relações com os conteúdos abordados em cada uma delas.
Os comentários da Banca de docentes corretores, a respeito do observado na correção das respostas dos estudantes, suas impressões e conclusões serão apresentados junto à análise de cada questão.
Cumpre esclarecer que, tendo -se em vista que as questões discursivas de Formação Geral são padronizadas, ou seja, constam de todas as provas, os comentários da Banca são os mesmos para todas as carreiras acadêmicas, sendo direcionados a todos os estudantes que participaram do Enade/2018.
A seguir, serão analisados os desempenhos da Área de Administração nas duas questões discursivas de Formação Geral do Enade/2018, comparando -se os resultados obtidos com comentários para cada questão. 6.3.1.1 Análise de Conteúdo da Questão Discursiva 1 do Componente de Formação Geral Os dados de Administração, obtidos a partir das respostas à questão 1, encontram -se na Tabela 6.14 e no Gráfico 6.7.
Nessa questão – de desempenho muito semelhante ao da outra questão de Formação Geral – os alunos de todo o Brasil, tiveram Média 32,5.
A maior Média para a questão 1 foi obtida na região Norte ( 34,5 ), e a menor, na região Centro-Oeste ( 30,4 ).
Quanto à variabilidade das notas, o Desvio padrão de todo o Brasil foi 21,7.
O menor Desvio padrão foi obtido na região Sul ( 21,3 ), e o maior Desvio padrão foi obtido pelas regiões Norte e Nordeste ( 22,0 ).
A Mediana do Brasil, como um todo, foi 30,0, sendo igual em quase todas as regiões, com exceção da região Nordeste ( 40,0 ).
As notas Máxima e Mínima da questão discursiva 1 foram, sem exceção, as mesmas para todas as regiões do Brasil, respectivamente, 100,0 e zero.
No Gráfico 6.7, mostra -se a distribuição das notas na questão discursiva 1 do Componente de Formação Geral.
Observa -se que quase 10 \% dos participantes deixaram em branco a resposta a essa questão.
Somados à frequência dos que, tendo respondido, receberam nota zero, chega próximo de 17 \% o contingente daqueles que zeraram questão.
O intervalo modal é ( 40,50 ] que caracteriza o máximo local com, aproximadamente, 20 \%, seguido de perto do intervalo ( 30 ; 40 ]. 6.3.1.2 Comentários sobre a correção de Conteúdo das respostas à Questão Discursiva 1 A Questão 1 exigia do aluno um posicionamento sobre conteúdo humanista e crítico.
O comando da questão solicitava a redação de um texto que formulasse e articulasse argumentos consistentes envolvendo temas previstos na Portaria nº 444, de 30 de maio de 2018, que dispõe sobre o componente de Formação geral do Enade 2018, tais como : solidariedade/violência, tolerância/intolerância, inclusão/exclusão, sexualidade, relações de gênero e relações \textbf{étnico-raciais} ( art. 7º, item VIII ), bem como ética, democracia e cidadania ( art. 7º, item I ).
O enunciado da questão buscava o estabelecimento da relação de três dimensões – os Direitos Humanos, a Democracia e as Fake News.
Para isso, recorreu a quatro suportes : a reprodução de uma foto e três fragmentos adaptados de textos que foram disponibilizados em páginas da Internet ( da Anistia Internacional, do Jornal El País-Brasil e do Labic-Laboratório de estudos sobre Imagem e Cibercultura ).
Os fragmentos traziam dados sobre o aumento do número de assassinatos de grupos minoritários, as origens e o trabalho da vereadora Marielle Franco, sua luta pelos Direitos Humanos, seu assassinato, e a produção de fake news. vereadora Marielle Franco, sua luta pelos Direitos Humanos, seu assassinato, e a produção de fake news.
O comando solicitava um texto abordando dois tópicos : o tensionamento entre a defesa dos Direitos Humanos e a produção de notícias falsas, tendo como base a vereadora Marielle – a defesa dos Direitos Humanos realizada por ela e as notícias falsas divulgadas após a sua morte.
As fake news também estavam presentes no segundo tópico de análise : os prejuízos que provocam para a sociedade democrática.
Por um lado, os textos selecionados para constituir o enunciado da questão continham apoio suficiente para formulação e articulação de argumentos para discorrer sobre o tensionamento entre a atuação da Marielle em defesa dos Direitos Humanos e fake news.
Por outro, havia ênfase nos aspectos biográficos e na trajetória da vereadora Marielle, o que levou a respostas pouco focadas no comando, que discutiam a vida e a morte da vereadora, sua importância como um ícone dos Direitos Humanos, opiniões em torno do seu assassinato, até mesmo sobre a responsabilidade pelo crime, muitos textos com a presença/defesa de fake news ou produção de novas.
O foco sobre Marielle também ensejou respostas basicamente \textbf{centradas} na questão da violência.
Foram frequentes as interpretações com viés político-partidário, não previstas no padrão de resposta.
O item do comando que solicitava a discussão sobre democracia e fake news encontrou menos apoio nos textos.
Apesar de se esperar que graduandos tivessem maturidade, inclusive política, para demonstrar alguma compreensão dos aspectos principais que envolvem o convívio democrático, esse aspecto foi pouco discutido e, quando presente, quase sempre foi apenas mencionado, numa quase reprodução do comando.
A falta de entendimento do que é o conceito de Direitos Humanos também foi percebida nas respostas, seja para recomendá -lo como importante para a sociedade, seja para criticá -lo como algo que defende bandidos.
Foi marcante a ideia que vincula Direitos Humanos a direitos apenas das minorias.
Não foi incomum encontrar respostas que \textbf{personificavam} os Direitos Humanos, tratando -os como sujeito de alguma ação que deveria ser implementada na sociedade em geral ou em casos específicos.
As fake news foram citadas muitas vezes, mas pouco explicadas, conforme previsto no padrão de resposta.
O argumento utilizado geralmente vinculava as consequências de fake news para a imagem da Marielle, ou seja, centrava no fato em si, já indicado no enunciado, sem qualquer análise ou aprofundamento próprio.
É importante salientar que muitas respostas, provavelmente de grupos vinculados a cursos específicos, continham argumentações recorrentes : os Direitos Humanos como algo vinculado à Constituição e a tratados dos quais o Brasil é signatário ; as fake news como produto do mundo digital e globalizado, relacionado à era da pós-verdade ; a discussão sobre o papel das mídias sociais e da imprensa.
Como já apontado, houve baixa quantidade de respostas contemplando o segundo tópico do comando – prejuízos da produção de notícias falsas para a sociedade democrática ; fato motivado, possivelmente, pela falta de um elemento que ressaltasse esse ponto nos textos motivadores.
De modo geral, os estudantes não conseguiram recorrer a conhecimentos tácitos sobre democracia para construir argumentos próprios.
As respostas que abordavam o tema, geralmente estavam vinculadas à ideia do direito à verdade, do direito à integridade da pessoa ou dos direitos gerais garantidos pela Constituição.
A relação da questão com o momento político brasileiro foi indicada em um número expressivo de respostas.
Por vezes, apenas questionando a temática da questão – caso Marielle – em uma prova para avaliar o ensino superior no Brasil ou a importância dos Direitos Humanos na atual conjuntura política do país, levando à anulação da resposta.
Houve também interpretações das consequências do uso das fake news no processo eleitoral de 2018, influenciando a escolha consciente e livre das pessoas.
As respostas demonstraram atenção político social frente a esses temas, porém com pouco domínio conceitual do campo de conhecimento sobre Direitos Humanos, fake news e democracia.
Em termos da linguagem adotada pelos alunos, há dois grandes grupos de respostas.
No primeiro grupo, os que responderam de forma clara, com autonomia de pensamento e explicitação de justificativas para os aspectos solicitados no enunciado, demonstrando conhecimento sobre o que significam Direitos Humanos, fake news e democracia, além da articulação entre essas questões.
Por outro lado, um segundo grupo dos que não conseguiram ir além do que foi citado no enunciado, apenas repetindo -o, muitas vezes desviando o foco da questão para temas correlatos, como a vida de Marielle ou a violência.
Nesse segundo grupo, mais facilmente foram encontradas respostas desconexas e mal articuladas, com textos confusos que explicitam o não conhecimento sobre o tema e a falta de adequação à linguagem, necessária a um concluinte do ensino superior.
As abordagens dos Direitos Humanos mais recorrentes nas respostas corretas foram : a referência da Marielle como representante/ícone dos DH ; e os DH como direitos previstos na Constituição e em tratados internacionais dos quais o Brasil é signatário.
No entanto, também foram encontradas críticas aos Direitos Humanos, tais como : DH vistos como “ direitos de bandidos ” ; a defesa pouco embasada de frases de efeito como “ direitos humanos para humanos direitos ”.
Também chamou atenção o tratamento dos DH como um personagem, ou um órgão público, evidenciando desconhecimento sobre o tema.
Em relação às fake news, dois argumentos previstos pelo padrão de respostas foram os mais presentes : a noção de que o mundo globalizado e a facilidade de reprodução de informações pelas mídias atuais levam à criação de \textbf{mentiras} que são compartilhadas e que geram consequências ; a Marielle como alvo de fake news que tentaram desconstruir sua luta e seu legado.
Também foi frequente a ideia de que a produção de informações falsas é realizada por uma imprensa tendenciosa com intenção de manipular a informação.
No item relacionado a fake news e democracia, as respostas corretas afirmavam que as fake news inviabilizam o acesso à verdade e à garantia da integridade de um Estado Democrático de Direito.
Também foi abordada a influência das fake news nos resultados das eleições dos últimos anos, nos Estados Unidos e no Brasil, gerando consequências para a vida das pessoas.
Por fim, um outro aspecto abordado nas respostas a este tópico foi a polarização da sociedade em termos políticos, com a afirmação de que as fake news dificultam o debate, estimulam posições extremas e rivalidades prejudiciais à convivência em uma democracia plena.
A notas mais fracas ( de zero a 35 ) foram cerca de 20 \% das respostas corrigidas – as que não estavam em branco nem foram desconsideradas ou anuladas.
Foram respostas que se restringiam a cópia ou paráfrases do enunciado ; as que tinham foco em desdobramentos políticos e policiais do caso Marielle ; as que tinham argumentação confusa e/ou incompleta.
A grande maioria das notas, em torno de 65 \%, foi mediana, entre 40 a 70.
Nesse grupo ficaram as respostas que continham alguma abordagem adequada sobre Direitos Humanos, sobre as atividades da vereadora Marielle e sobre as fake news divulgadas após sua morte.
As notas mais altas deste grupo foram para respostas que incluíram a questão do tensionamento, quase sempre, em um texto coerente e consistente.
Muitas vezes, a nota mediana foi consequência do baixo índice de respostas à segunda parte do comando, que solicitava a discussão sobre os prejuízos das fake news para a democracia.
As melhores notas só foram obtidas por aqueles que incluíram a discussão sobre os prejuízos das fake news para a democracia.
Apenas 15 \% das respostas corrigidas receberam avaliação igual ou superior a 75.
A deficiência principal detectada pelos corretores foi a dificuldade de grande parte dos alunos para interpretar o enunciado da questão e para se expressar em textos consistentes.
De modo geral, os temas foram tratados de forma fragmentada e pouco analítica.
Além disso, chamou à atenção a predominância do lugar comum na construção argumentativa das respostas, a superficialidade do debate em torno da democracia e o desconhecimento de conceitos como Direitos Humanos, esses confundidos com instituições e até pessoas. 6.3.1.3 Análise de Conteúdo da Questão Discursiva 2 do Componente de Formação Geral A Tabela 6.15 mostra que o desempenho médio dos estudantes, na questão discursiva 2 ( média 32,6 ), foi praticamente igual ao obtido na questão discursiva 1 ( média 32,5 ).
A região Norte foi aquela cuja Média, também nessa questão, foi maior ( 34,8 ), assim como a de menor Média foi a região Centro-Oeste ( 30,4 ).
Quanto à variabilidade das notas, o Desvio padrão de todo o Brasil foi 21,4, um pouco inferior ao obtido na questão discursiva 1 ( 21,7 ).
O maior desvio nessa questão foi obtido na região Norte ( 22,0 ), enquanto o menor foi obtido na região Sul ( 20,5 ).
A Mediana de todo o Brasil foi 35,0, a mesma de três regiões.
As exceções foram as regiões Sul e Centro-Oeste, que obtiveram Mediana 30,0.
As notas : Máxima ( 100,0 ) e Mínima ( 0,0 ), foram as mesmas em todas as regiões, sem exceção.
No Gráfico 6.8, mostra -se a distribuição das notas na questão discursiva 2 do Componente de Formação Geral.
Observa -se que, um pouco mais de 10 \% dos alunos deixaram a questão em branco.
Chega a 15 \% se agregado à frequência daqueles que, tendo respondido, receberam nota zero.
O intervalo de maior frequência foi ( 30 ; 40 ], seguido do intervalo ( 40 ; 50 ].
Nota -se, ainda, que as notas ficaram levemente menos dispersas em comparação às da questão discursiva de número 1, o que pode ser constatado, também, pela comparação do desvio padrão das notas da questão discursiva 2 ( 21,4 ) e o das notas da questão discursiva 1 ( 21,7 ). 6.3.1.4 Comentários sobre a correção de Conteúdo das respostas à Questão Discursiva 2 Os textos e o comando da Questão 2 de Formação Geral são perfeitamente adequados ao que se espera avaliar em Formação Geral, ao tratar das funções sociais dos museus, contextualizando a temática sob o ponto de vista da memória e da perspectiva de futuro.
Na questão, estavam presentes, principalmente, o conteúdo humanista e crítico, a interpretação e a produção de textos sobre temas como cultura e arte, educação, ciência e democracia, previstos na Portaria que dispõe sobre o componente de Formação Geral do Enade 2018.
Os conhecimentos exigidos para responder ao comando, explicitados no padrão de resposta, vão além do senso comum, por se exigir a compreensão de conceitos como “ função social ” ( de instituições sociais ), “ memória ” ( no sentido de memória social ), “ perspectiva de futuro ” ( no contexto \textbf{sócio-histórico-cultural} ) e “ sociedade contemporânea ” ( que deveria independer da interpretação de cada grupo social a respeito da sociedade na qual vive ).
No entanto, apesar de envolver tais conceitos, espera -se que alunos concluintes de cursos superiores de graduação sejam capazes de discutir funções de instituições sociais, especialmente porque o enunciado trazia fato recente e de grande repercussão como apoio.
O padrão de resposta elencou as funções sociais que poderiam ser abordadas nas respostas dos estudantes, como também argumentos que poderiam ser usados ao longo das mesmas.
A tarefa dos estudantes era a construção de um texto expositivo que articulasse a importância de museus, suas funções sociais e as perspectivas de memória e de futuro, o que se mostrou de dificuldade de média a difícil.
Apesar de o enunciado ser claro, muitos estudantes, provavelmente influenciados pelos textos, limitaram -se a discutir aspectos relacionados ao Museu Nacional e, em particular, ao incêndio que o \textbf{destruiu}.
Foram poucos os que interpretaram corretamente o escopo do comando, que solicitava a discussão das funções sociais de museus de uma forma geral, no sentido de argumentar a importância deste tipo de instituição.
As quatro funções sociais elencadas no padrão de resposta foram encontradas na quase totalidade das redações dos estudantes de forma articulada com a importância dos museus.
Entretanto, quase nunca os textos abordavam mais do que uma função social e vinculavam, de forma clara, a importância desse tipo de instituição para a sociedade contemporânea.
Além disso, a solicitação de contextualização sob o foco museu/memória/perspectiva de futuro parece não ter sido bem compreendida e, em consequência, esteve pouco presente nas respostas.
Também se observou confusão entre os conceitos de história e memória e uma visão estreita do significado de sociedade contemporânea, às vezes, restrita à cidade do respondente, outras, ao Rio de Janeiro, ou ainda a um tipo de sociedade idealizada.
A relevância que os textos de apoio deram ao caso particular do incêndio do Museu Nacional do Rio de Janeiro e aos prejuízos causados ao seu acervo, induziu a maioria dos estudantes a elaborar respostas tratando das funções a partir da perda e, em particular, à perda do tipo de acervo daquele museu e do tipo de pesquisa que lá se desenvolve, divergindo do comando da questão, que solicitava texto abordando funções sociais de museus em geral.
Também foram encontradas muitas respostas que se limitavam a discutir o que o estudante considerava como possíveis causas do incêndio, a falta de manutenção, questões políticas e de gestão, sem responder ao que foi solicitado.
Foi também recorrente a transcrição de frases dos textos do enunciado nas respostas, diretamente ou por paráfrases, revelando certa dificuldade de criação textual.
A análise das respostas dos estudantes, além de revelar dificuldade de abordar a temática esperada, revelou outros dois aspectos predominantes.
Um deles, como já comentado, foi a ênfase no incêndio do Museu Nacional, que orientou grande parte dos textos dos estudantes.
Outro, também evidenciando desvio do esperado, foi o foco na perda patrimonial.
Nesse grupo, houve respostas que até apresentaram uma função social, mas olhada pelo ponto de vista da perda.
Nessa linha, muitas respostas enalteciam a importância do Museu Nacional, sua relevância no cenário nacional, o impacto da perda de seu acervo, isso complementado por frágeis reflexões sobre as funções sociais.
Em consequência, a perspectiva de desdobramentos mais ricos, que tratassem de aspectos sociais relacionados a diferentes tipos de museu foi prejudicada.
Muitas respostas trataram da preservação da memória e/ou guarda do acervo ; identidade, tanto individual quanto coletiva ; conhecimento sobre antepassados ; evolução da sociedade e dos seus costumes ; valorização e disseminação do conhecimento por meio da liberação do acesso aos acervos ; bem como atividades pedagógicas e exposições de arte.
O aspecto produção de pesquisa e conhecimento, embora frequente, foi mencionado de forma muito simples, sem qualquer aprofundamento, para além do que já era mencionado no enunciado.
Frequentemente apresentava -se o museu como produtor de pesquisas para o campo da saúde, o que evidencia desconhecimento das reais possibilidades de pesquisa, mesmo no Museu Nacional.
Quanto a perspectiva de futuro, quando presente no texto, era mencionada de forma breve, com afirmações genéricas sobre a possibilidade de melhoria do futuro pelo simples conhecimento do passado, numa percepção de que o passado determina o futuro, de forma linear.
Foram raras as respostas com alguma argumentação que associasse pesquisa, memória, identidade e reflexão social como base para a construção de alternativas de futuro.
A relação entre turismo e museu apareceu considerando a visita a museus como forma de lazer e, algumas vezes, como geradora de recursos.
O desempenho pode ser considerado fraco.
Muitas respostas deixavam claro que o estudante nunca tinha visitado um museu e, de uma forma geral, observou -se desconhecimento sobre suas funções sociais.
Não se esperava que história, memória social e museologia fossem assuntos dominados pelos estudantes das diversas áreas que participaram do Enade 2018, mas esperava -se que essa temática pudesse ser abordada do ponto de vista das experiências pessoais ou de conhecimento social.
Como uma vivência própria com os museus também se revelou frágil ou inexistente, o conhecimento sobre o assunto ficou bem restrito aos textos motivadores e ao senso comum, construído, basicamente, pelo noticiário em torno do incêndio.
O universo dos argumentos dos estudantes se mostrou muito limitado.
Cada dimensão foi praticamente representada por uma ideia simplista.
De forma recorrente : a preservação da memória para construção da identidade de um povo ; a compreensão do passado por meio das pesquisas como oportunidade para melhorar o futuro ; o museu como um espaço em que há um \textbf{aprendizado} estimulado pelo lazer.
Houve, também, uma grande concentração de respostas associando o museu a um local de guarda de acervo e de história.
As respostas à Questão 2 evidenciaram dificuldades conceituais.
Além dos conceitos de história, memória e museu se misturarem, parecendo não se distinguirem, ao mencionarem a pesquisa, também se observou falta de distinção entre pesquisas colegiais, com objetivos de \textbf{aprendizado}, e as realizadas por pesquisadores, com objetivos de avançar cientificamente.
A falta de compreensão do que se pedia no comando da questão e a dificuldade de expressão linguística foram dois aspectos marcantes na avaliação dos corretores.
Foi frequente os textos apresentarem ideias soltas ou de forma confusa, sem constituir unidade.
O vocabulário usado se revelou bem restrito, e fortemente apoiado no coloquial, na linguagem oral.
Predominou o senso comum em lugar de uma reflexão mais formal envolvendo os conceitos envolvidos. 6.3.1.5 Análise de Língua Portuguesa das Questões Discursivas do Componente de Formação Geral Os dados de Administração, obtidos a partir das respostas às questões discursivas do Componente de Formação Geral, no que tange à Língua Portuguesa, encontram -se na Tabela 6.16 e no Gráfico 6.9.
Nesse aspecto, os alunos de todo o Brasil obtiveram Média 55,3.
A maior Média em relação à Língua Portuguesa foi obtida na região Sudeste ( 56,1 ), e a menor, na região Centro-Oeste ( 52,4 ).
Quanto à variabilidade das notas, o Desvio padrão de todo o Brasil foi 21,4.
O menor Desvio padrão foi obtido na região Sul ( 20,7 ) e o maior Desvio padrão foi obtido na região Centro-Oeste ( 22,9 ).
A Mediana das notas de Língua Portuguesa foi 62,5 para o Brasil, como um todo, a mesma obtida nas regiões Norte e Sudeste.
A Mediana das regiões Nordeste, Sul e Centro-Oeste foi 60,0.
A nota Máxima para todo o Brasil foi 92,5, com, pelo menos, um aluno tirando essa nota em quase todas as regiões, exceto na região Centro-Oeste ( 90,0 ).
Já a nota Mínima foi zero em todas as regiões do país. das pesquisas como oportunidade para melhorar o futuro ; o museu como um espaço em que há um \textbf{aprendizado} estimulado pelo lazer.
Houve, também, uma grande concentração de respostas associando o museu a um local de guarda de acervo e de história.
As respostas à Questão 2 evidenciaram dificuldades conceituais.
Além dos conceitos de história, memória e museu se misturarem, parecendo não se distinguirem, ao mencionarem a pesquisa, também se observou falta de distinção entre pesquisas colegiais, com objetivos de \textbf{aprendizado}, e as realizadas por pesquisadores, com objetivos de avançar cientificamente.
A falta de compreensão do que se pedia no comando da questão e a dificuldade de expressão linguística foram dois aspectos marcantes na avaliação dos corretores.
Foi frequente os textos apresentarem ideias soltas ou de forma confusa, sem constituir unidade.
O vocabulário usado se revelou bem restrito, e fortemente apoiado no coloquial, na linguagem oral.
Predominou o senso comum em lugar de uma reflexão mais formal envolvendo os conceitos envolvidos.
No Gráfico 6.9, mostra -se a distribuição das notas de Língua Portuguesa do Componente de Formação Geral.
Observa -se que a maior frequência ( um pouco mais do que 30 \% ) corresponde à dos alunos que obtiveram nota no intervalo ( 60 ; 70 ].
Destacam -se, também, os alunos que deixaram ambas as questões em branco, representando, aproximadamente, 7 \% do total. 6.3.1.6 Comentários sobre a correção das respostas de Formação Geral com respeito à Língua Portuguesa As duas questões discursivas de Formação Geral ensejaram duas configurações textuais distintas, tendo em vista os seus encaminhamentos : a Questão 1 teve um caráter mais opinativo e a Questão 2, mais expositivo, conforme se explicita a seguir.
Ambas as questões permitiram a \textbf{observação} do desempenho linguístico dos participantes em situação formal, competência exigida no processo de formação universitária.
As habilidades de produção textual envolvidas têm, igualmente, como condição básica, o domínio do padrão formal culto da Língua Portuguesa, no que diz respeito aos aspectos relativos às questões de caráter ortográfico, textual, morfossintático e vocabular.
O encaminhamento das questões evidencia a expectativa de que o participante utilizasse seus conhecimentos sobre os dois temas e estruturasse seus textos de acordo com as características do registro formal adequado à situação comunicativa.
Tal configuração determina exigências quanto : à adequação da seleção vocabular, ao desenvolvimento do conteúdo, à estruturação sintática dos períodos, à organização lógica das ideias, à utilização de procedimentos de encadeamento textual e de referenciação, à obediência às exigências morfossintáticas próprias da modalidade escrita da norma-padrão, ao respeito às regras ortográficas e às regras de acentuação gráfica.
Esse encaminhamento exigiu do graduando o domínio do modo de organização textual de base dissertativo-argumentativa, essencial ao processo de formação universitária.
As habilidades de produção textual envolvidas têm como condição básica o domínio do padrão formal culto da Língua Portuguesa, no que diz respeito aos aspectos relativos às questões de caráter ortográfico, textual, morfossintático e vocabular.
O enunciado da Questão 1 tomava como base três fragmentos de textos jornalísticos publicados em sites da internet, que abordavam o aumento do assassinato de minorias e de defensores de direitos humanos no Brasil, entre eles a vereadora Marielle Franco.
O comando solicitava que o participante desenvolvesse um texto dissertativo e argumentativo com as temáticas Fake News, Democracia e Direitos Humanos.
Embora não tenha havido referência explícita à tipologia textual esperada para o desenvolvimento da resposta, essa proposta encaminhou o participante para a elaboração de um texto opinativo, com o objetivo de discutir o problema e manifestar um posicionamento crítico.
O enunciado da Questão 2 abordava o incêndio do Museu Nacional do Rio de Janeiro, ocorrido em 2018, e seus impactos na ciência brasileira e internacional, tomando como base dois fragmentos de textos jornalísticos publicados em sites da internet.
O comando solicitava que o participante desenvolvesse um texto abordando três aspectos da função social dos museus na sociedade contemporânea, encaminhando -o para a elaboração de um texto expositivo que confirmasse a posição defendida nos fragmentos motivadores a respeito da importância dessas instituições.
Para atender a essas exigências, o padrão de resposta utilizado na avaliação das questões discursivas de Formação Geral considerou aspectos relevantes ao bom desempenho linguístico como competências distintas, de modo a permitir um mapeamento detalhado do domínio dos recursos disponíveis na Língua Portuguesa para a comunicação escrita formal : ortográficos, textuais, morfossintáticos e vocabulares ( ver detalhes no Anexo VIII – Padrão de Respostas ).
Aspectos Ortográficos Para avaliar a competência relativa ao domínio das convenções ortográficas da norma-padrão da Língua Portuguesa, observou -se o respeito às regras de acentuação gráfica e da grafia padrão das palavras ( com ausência de abreviaturas próprias da linguagem da internet ), de acordo com as convenções estabelecidas pela legislação em vigor e consubstanciadas no Vocabulário Ortográfico da Língua Portuguesa, editado pela Academia Brasileira de Letras.
Aspectos Textuais Para avaliar a competência relativa ao domínio dos procedimentos de estruturação textual, procurou -se observar a utilização de mecanismos linguístico-discursivos que estabeleçam o encadeamento lógico entre as partes do texto, de forma a garantir a progressão e a coerência textuais.
Esta competência destaca -se por sua relevância, já que são avaliadas a clareza e a unidade textual.
Aspectos Morfossintáticos Para avaliar o domínio dos diferentes aspectos morfossintáticos próprios da modalidade escrita formal da norma-padrão da Língua Portuguesa, observou -se se o texto produzido atendeu às seguintes exigências : a concordância nominal, a concordância verbal, a regência nominal, a regência verbal, a flexão nominal, a flexão verbal, a correlação entre os tempos verbais, a colocação pronominal e a utilização de sinais de pontuação que contribuam para a organização lógica da frase e do texto.
Aspectos Vocabulares Para avaliar a adequação da seleção vocabular à modalidade escrita formal da Língua Portuguesa, exigida pela situação comunicativa, observou -se o atendimento às seguintes exigências : precisão na seleção/utilização do vocabulário relacionado à temática solicitada pela questão ; ausência de marcas da oralidade, como termos de sentido muito genérico ( “ coisa ”, “ negócio ”, “ você ” ) e termos de registros mais informais ( como gírias, jargões, frases feitas, ditados populares, termos regionais ).
Assim, espera -se que o participante respeite a adequação vocabular não usando gírias ou expressões coloquiais, evite repetição desnecessária de palavras e utilize um vocabulário mais formal, como solicitado por um texto dissertativo. ( “ coisa ”, “ negócio ”, “ você ” ) e termos de registros mais informais ( como gírias, jargões, frases feitas, ditados populares, termos regionais ).
Assim, espera -se que o participante respeite a adequação vocabular não usando gírias ou expressões coloquiais, evite repetição desnecessária de palavras e utilize um vocabulário mais formal, como solicitado por um texto dissertativo.
Para efeito de pontuação, a avaliação do desempenho linguístico considerou três grandes grupos de competências, agrupando aspectos explicitados, com os seguintes \textbf{pesos} relativos : aspectos ortográficos ( 20 \% ) ; aspectos textuais ( 40 \% ), aspectos morfossintáticos e vocabulares ( 40 \% ).
Os critérios de pontuação utilizados no padrão de resposta procuraram atender à diversidade de desempenho dos graduandos.
Nos textos analisados, observaram -se, assim, vários níveis de atendimento às exigências do padrão formal próprio da modalidade escrita da Língua Portuguesa, configurando um"continuum".
Por um lado, textos que traduzem com exatidão a expectativa de uso formal da língua, em todos os níveis analisados, marcados por clareza e unidade ; e, no outro extremo, textos quase incompreensíveis devido ao grande número de desvios de todos os tipos, principalmente os de caráter textual.
As duas questões discursivas de Formação Geral analisadas revelaram desempenhos distintos dos participantes.
A Questão 1 propiciou a oportunidade de maior desenvolvimento do tema solicitado, gerando consequências na estruturação textual, já que os textos foram mais longos, mais elaborados e fluentes.
Em virtude da polêmica sobre a existência das fake news e sobre a visão pública da personagem da vereadora Marielle Franco, observaram -se respostas que polemizaram o tema, favorecendo a construção argumentativa do texto.
A Questão 2 teve um comportamento distinto, devido ao caráter mais expositivo do enunciado e ao reduzido conhecimento dos participantes sobre o tema.
Assim, os textos foram mais curtos, frequentemente com repetição de ideias dos textos motivadores.
Quanto aos aspectos linguísticos analisados durante esta avaliação, observaram -se os seguintes resultados : Aspectos ortográficos : O desempenho dos participantes revelou uma diferença muito grande nos dois aspectos analisados nesta competência : baixo índice de desvios da grafia padrão e grande índice de desvios de acentuação.
Em vários casos, ocorre ausência completa de acentuação gráfica em todas as palavras do texto.
Os resultados revelaram que a tendência dominante entre os universitários brasileiros é a eliminação da acentuação gráfica, talvez motivada pelos hábitos relacionados às redes sociais e pela ausência de esclarecimento dos meios de comunicação, das autoridades e das escolas sobre as decisões do Acordo Ortográfico da Língua Portuguesa de 1990.
Os casos mais sistemáticos de eliminação do acento indicador da sílaba tônica foram : - palavras proparoxítonas ( por exemplo “ generos ”, “ arbitro ”, “ politico ”, “ publicas ”, “ numero ” ) ; - palavras paroxítonas terminadas em ditongo crescente ( por exemplo “ area ”, “ noticiarios ”, “ individuo ”, “ historia ”, “ varios ”, “ noticia ”, “ negligencia ”, “ ocorrencia ” ) ; - palavras oxítonas ( por exemplo “ ninguem ”, “ esta ”, “ ate ”, “ tambem ” ).
Por outro lado, destaca -se o uso indevido do acento gráfico em palavras como “ melâncolia ” e “ intervênção ”, por exemplo.
Quanto ao domínio das convenções relativas à grafia das palavras, observaram -se desvios \textbf{esporádicos} como : hipercorreção pela escolha de “ e ” no lugar de “ i ”, por influência de hábitos da oralidade ( “ descriminação ” no lugar de “ discriminação ”, “ entervenção ” no lugar de “ intervenção ” ) ; desvios de grafia relacionados à variação diastrática, como por exemplo “ perca ”, “ subjulgadas ”, “ precoseito ”, “ soubre ”, “ vecendo ”, “ indesencia ”, “ apolojia ”, “ fulga ”, “ dereitos ” ;...
Observaram -se, também, muitos casos de inadequação no uso da maiúscula : uso indevido para destacar determinadas palavras-chave do texto, como “ Brasileiros ”, “ Homens ” ; ausência de maiúscula para grafar nomes próprios, como"marielle".
Expressões como"Direitos Humanos"foram grafadas diversificadamente, com maiúsculas e com minúsculas.
Destaque -se, também, a maioria dos participantes que grafam os textos inteiramente em caixa alta, o que impede a identificação de desvios relacionados ao uso de maiúsculas.
Vale observar, também, que, ao contrário do que se esperava, não apareceram abreviaturas próprias do “ internetês ”, relacionadas ao uso de redes sociais e e-mails, como por exemplo"vc".
Aspectos textuais : Para alguns dos graduandos essa competência se revelou como a mais problemática, tendo em vista os inúmeros problemas observados, desvios acumulados durante toda a formação escolar.
São eles : sequência justaposta de ideias sem encaixamentos sintáticos, provocando uma redução de estruturas subordinadas, ao lado do aumento na frequência de estruturas coordenadas e absolutas ; frases fragmentadas que comprometem a estrutura Aspectos textuais : Para alguns dos graduandos essa competência se revelou como a mais problemática, tendo em vista os inúmeros problemas observados, desvios acumulados durante toda a formação escolar.
São eles : sequência justaposta de ideias sem encaixamentos sintáticos, provocando uma redução de estruturas subordinadas, ao lado do aumento na frequência de estruturas coordenadas e absolutas ; frases fragmentadas que comprometem a estrutura lógico-gramatical ; frases formadas apenas por oração subordinada, sem oração principal ; redução no uso de conectores para expressar relações lógicas essenciais à construção do texto como consequência da mudança de estruturação frasal ; emprego inadequado do pronome relativo ( com omissão da preposição ou a utilização do pronome relativo inadequado, como “ onde ” ), refletindo hábitos da oralidade ; ausência de recursos de referenciação, como substituição de termos por sinônimos, hiperônimos, nominalizações, expressões metafóricas.
Esses problemas revelam dificuldades em relação à estrutura formal do texto produzido, o que é preocupante ao se levar em conta que são graduandos em fase final de formação.
Em alguns textos, falta um mínimo de textualidade e de domínio do registro padrão da língua.
Na verdade, observaram -se relações linguísticas quase agramaticais, como as estabelecidas pela sequência de gerúndios sem o apoio de um ponto de partida para a organização das informações gramaticais e semânticas.
Quanto à utilização dos sinais de pontuação, observou -se uma ausência dos mesmos na maioria dos textos analisados.
Foi muito frequente a ocorrência de parágrafos sem marca interna de pontuação para separar os períodos.
São os seguintes os tipos de problemas mais encontrados :  vírgula : utilização de vírgula para separar o sujeito e o predicado ; uso de vírgula no lugar do ponto para separar ideias que constituem períodos distintos ; ausência de vírgula para separar oração adjetiva explicativa ;  ponto e vírgula : utilização do ponto e vírgula no lugar de vírgula ;  ponto final : ausência de ponto final para separar períodos.
Aspectos morfossintáticos e vocabulares : Os resultados são transparentes em relação aos aspectos mais problemáticos no desempenho dos participantes nestas duas competências.
Em relação aos aspectos morfossintáticos, seguem algumas \textbf{observações} sobre os desvios mais frequentes.
Quanto à regência, o desvio mais frequente foi a falta do sinal indicativo da crase – isso revela que o usuário não tem consciência de que, sob a forma do termo “ a ”, existe a presença de uma contração entre a preposição “ a ” ( exigida pela regência do termo anterior ) e o artigo definido “ a ”.
Outro problema relacionado à regência verbal e à nominal, encontrado frequentemente nas respostas, foi a ausência de preposição antes de pronome relativo, processo generalizado na modalidade oral da língua, em situações de registro informal.
Apesar da possibilidade de que essa alteração de regência se generalize no padrão escrito da Língua Portuguesa, como já está ocorrendo até em textos jornalísticos, o não emprego da preposição foi considerado como desvio neste processo de avaliação.
Outro desvio de regência significativo foi a utilização inadequada de uma preposição ou sua ausência após o verbo ou o nome ( substantivo ou adjetivo ).
A concordância verbal e a concordância nominal apresentaram alguns desvios muito frequentes.
Quanto à concordância de número, observou -se, como apontam as pesquisas nessa área, ausência de marca ( com sujeito anteposto ou posposto ) ou uso indevido ( uso inadequado da marca de plural comandado pelo núcleo plural da locução adjetiva, apesar de o substantivo que funciona como núcleo do sintagma nominal estar no singular ).
Uma ocorrência que se destacou foi a ausência de acento circunflexo na forma plural do presente do indicativo do verbo “ ter ”, que foi considerada como um desvio na concordância verbal e não na acentuação gráfica.
Quanto à concordância de gênero, vários casos foram observados, normalmente no âmbito de sintagmas nominais longos, em que o adjetivo está afastado do substantivo.
Quanto à questão da colocação pronominal, foram poucos os desvios observados.
Concluiu -se que, no registro escrito formal, a maioria dos participantes já incorporou regras como a não introdução da frase por um pronome oblíquo e a próclise na presença de um termo atrator.
Várias marcas de oralidade foram identificadas, embora não com alta frequência : uso do pronome relativo “ onde ” como relativo universal, falta de artigo definido antes de substantivo isolado ou antes de termos em enumerações.
Quanto aos aspectos vocabulares, alguns tipos de inadequação foram observados : expressões da oralidade ; seleção vocabular incompatível com o contexto, gerando falta de inteligibilidade ; falta de domínio de vocabulário mais abstrato e de maior complexidade, essencial ao desenvolvimento do texto de base dissertativa.
O principal aspecto observado foi a excessiva repetição de certas palavras, como o termo"pessoa", por exemplo, revelando limitação de repertório vocabular.
Em função do tema solicitado na Questão 1, diferentes grafias da expressão fake news não foram consideradas.
Para sintetizar, os problemas que mais se destacaram nos textos analisados foram : a ) truncamentos sintáticos, caracterizados por estruturas frasais incompreensíveis devido à ruptura da complexidade sintática própria da modalidade escrita ; b ) redução de estruturas subordinadas, compensada pelo aumento na frequência de estruturas coordenadas e absolutas, por um lado, ou pela elaboração de estruturas textuais prejudicadas pelo excesso de ideias sem a devida conexão subordinativa ; c ) redução no uso de conectores para expressar relações lógicas essenciais à construção do texto, substituídos pela exigência de inferência por parte do interlocutor para suprir a sua ausência ; d ) redução no uso do subjuntivo, ao lado da ampliação do uso do indicativo, combinado a estruturas frasais coordenadas ou absolutas ; d ) empobrecimento do processo de referenciação, com a repetição de palavras, sem processos mais sofisticados de substituição ; e ) simplificação da marcação da categoria tempo na morfologia verbal ; f ) desrespeito às regras de concordância verbal e nominal próprias do padrão formal ; g ) ausência do acento grave indicador da crase, revelando falta de conhecimento dos critérios de regência verbal e nominal ; h ) falta de domínio de vocabulário mais abstrato e de maior complexidade, essencial ao desenvolvimento do processo dissertativo ; i ) redução do emprego da acentuação gráfica.
Como resultado da correção, merece atenção a constatação de que a modalidade escrita tem apresentado a tendência a uma evidente simplificação, aproximando -se das características da modalidade oral da Língua Portuguesa.
No caso dos textos de base dissertativa ( expositiva ) e de base argumentativa, inscritos em um registro formal, a distância entre as duas modalidades é ainda maior, o que provoca desvios recorrentes em todos os quatro aspectos analisados : ortográficos, textuais, morfossintáticos e vocabulares.
Destaca -se, nessa avaliação, o comprometimento textual, por meio de estruturas fragmentadas e/ou truncadas, rompendo a complexidade sintática esperada no padrão formal. 6.3.2 Componente de Conhecimento Específico Na parte da prova relativa às questões discursivas do Componente de Conhecimento Específico ( Tabela 6.17 ), observa -se que a Média foi bem mais baixa do que para as questões discursivas do Componente de Formação Geral.
Enquanto no Componente de Formação Geral a Média para estudantes de Administração de todo o Brasil foi 37,1, na parte de Conhecimento Específico, a Média foi 19,0.
A maior Média deste componente foi obtida pelos estudantes da região Sudeste ( 20,1 ), e a menor, pelos da região Centro-Oeste ( 16,4 ).
As demais regiões obtiveram Média igual a 18,2.
Quanto à variabilidade das notas, o Desvio padrão de todo o Brasil foi 16,8.
O maior Desvio padrão foi encontrado na região Sudeste ( 17,3 ), e o menor, na região Centro-Oeste ( 15,7 ).
A maior nota Máxima, 100,0, foi obtida na região Nordeste.
Nas demais regiões, as notas máximas foram : 91,7, nas regiões Sudeste e Sul, 85,0, na Centro-Oeste e 83,3, na região Norte.
A nota Mínima ( 0,0 ) foi obtida por, pelo menos, um aluno em todas as regiões do Brasil.
A Mediana de todas as regiões foi a mesma do Brasil, como um todo, igual a 16,7.
No Gráfico 6.10, representa -se a distribuição das notas nas questões discursivas no Componente de Conhecimento Específico.
A moda dessa distribuição ocorre no intervalo [ 0 ; 10 ], com 38 \% do total de participantes.
Observa -se que nesses intervalos estão computadas as ocorrências de respostas em branco para as três questões discursivas de Conhecimentos Específicos.
A análise de cada uma destas questões será feita a seguir. 6.3.2.1 Análise da Questão Discursiva 3 do Componente de Conhecimento Específico Na questão 3, cujos resultados aferidos encontram -se descritos na Tabela 6.18, a Média dos estudantes de todo o Brasil foi 15,3.
O desempenho dos estudantes nessa questão foi superior ao da questão 4 e inferior ao da questão 5.
A menor Média nessa questão foi obtida pelos alunos da região Centro-Oeste ( 13,5 ), enquanto a maior Média foi obtida na região Sudeste ( 16,4 ).
Quanto à variabilidade das notas, o Desvio padrão de todo o Brasil foi 22,4.
O maior Desvio padrão foi obtido na região Sudeste ( 23,1 ), enquanto o menor foi obtido na região Centro-Oeste ( 20,7 ).
A nota Máxima, 100,0 pontos, foi alcançada por, pelo menos, um aluno em cada região.
A Mediana do Brasil, como um todo, foi zero, a mesma para todas as regiões.
A nota Mínima ( 0,0 ) também foi a mesma em todas as regiões do Brasil.
Cabe salientar que pelo menos a metade dos alunos em cada uma das regiões tirou nota zero nessa questão.
No Gráfico 6.11, mostra -se a distribuição das notas na questão discursiva 3, do Componente de Conhecimento Específico da área de Administração.
Sem se levar em consideração que 25,2 \% dos estudantes deixaram a questão em branco e que um pouco mais de 29 \% receberam nota zero, a moda dessa distribuição é o intervalo ( 20 ; 30 ], seguido do intervalo ( 0 ; 10 ]. 6.3.2.2 - Comentários sobre as respostas à Questão Discursiva 3 A questão apresentava conteúdo alinhado aos componentes da formação profissional do administrador, estando, portanto, afinada com o que é proposto nas Diretrizes Curriculares do curso.
Abarcava dois domínios importantes da formação, a saber, Teorias de Administração ( usualmente reunidas em Teoria Geral da Administração ) e Gestão da Produção, e fazia isso de forma correta e interessante, promovendo a interdisciplinaridade.
Dessa forma, no tocante aos conhecimentos exigidos do estudante, como também ao nível de profundidade, a proposta da questão era bem adequada.
Para responder correta e completamente à questão, o estudante deveria utilizar conhecimentos além dos óbvios ou que dependessem apenas de memorização.
Precisava ser capaz de relacionar o conhecimento de teorias às práticas de gestão e às áreas funcionais de uma organização.
A questão permitiu discriminação entre os diferentes níveis de conhecimento e habilidade dos estudantes para o estabelecimento dessas relações.
No tocante ao enunciado, a utilização de um texto introdutório, adaptado de um trecho de um livro de Peter Drucker, parece não ter contribuído para a elaboração de respostas aos comandos da questão.
Em apenas oito linhas são tratados diversos temas de administração, como avaliação de desempenho de empresas que visam ao lucro e de organizações sem fins lucrativos, desenvolvimento e sustentabilidade empresarial, gestão da produção e métodos da administração científica.
Além disso, o texto traz expressões de compreensão não trivial como “ nível de produção ” e “ linha de resultados ”, compondo a “ medição da administração e da empresa ”, enquanto, a rigor, não há uma maneira absoluta de medir a “ administração ” ou a “ empresa ”.
O texto introdutório também tinha foco específico na função de produção, ao mesmo tempo em que listava várias dimensões empresariais.
Isso pode ter provocado certa confusão de entendimento do comando, a qual foi observada nas respostas corrigidas.
A questão continha dois comandos.
O primeiro pedia ao respondente que relacionasse os métodos tayloristas a atividades de produção.
Já o segundo, completamente independente do primeiro, pedia que fossem deixados de lado aspectos de produção e operações, tendo o estudante que apresentar e explicar indicadores de desempenho gerencial para três áreas funcionais distintas.
Como resposta ao item ‘ a ’, muitas propostas poderiam ser consideradas como justificativas válidas, desde a associação dos métodos tayloristas ao aprimoramento de tarefas, até respostas simples como a associação entre Taylor e \textbf{chão} de fábrica, de forma genérica.
O comando do item ‘ b ’ era claro e direto.
Avaliava o conhecimento do estudante a respeito das áreas funcionais de uma organização, bem como avaliava se o estudante relacionava essas áreas à avaliação, conhecimento prático, fundamental para o futuro administrador.
Essa questão exigiu aplicação de conhecimentos de duas disciplinas distintas de forma ampla e bem estruturada e, portanto, poderia ser considerada a mais difícil entre as três discursivas de Conhecimentos Específicos.
O padrão de resposta foi apresentado em linguagem clara e acessível e com exigência de nível de profundidade adequado e propiciou uma correção justa, capaz de captar os diferentes níveis de respostas encontrados.
Com relação ao item ‘ b ’, o padrão listou algumas áreas que poderiam ser consideradas nas respostas dos estudantes – Finanças, Recursos Humanos, Marketing, Materiais e Logística, Sistemas de Informação – além de algumas delas apresentarem três exemplos de respostas completas.
No tocante à forma, percebeu -se alguma dificuldade dos estudantes na construção de textos coerentes e fluidos, especialmente nas respostas ao item ‘ a ’ que envolvia argumentação dissertativa para uma solicitação de apresentar uma justificativa.
As respostas apresentadas pelos estudantes ao item ‘ a ’ variaram, desde as raras completas e objetivas e muitas incompletas, que apenas relacionavam taylorismo e gestão da produção, sem apresentar justificativa.
Também se observaram muitas respostas erradas, que relacionavam Taylor à produção just in time ou ao toyotismo, confundiam Taylor e Elton Mayo, ou tratavam do experimento de Hawthorne.
Outras apenas repetiam o conteúdo do enunciado, ou traziam considerações genéricas sem oferecer nenhuma reflexão – algo como : “ Os métodos de Taylor aumentaram a lucratividade da empresa ”, ou “ o que Taylor propôs é utilizado até hoje ”.
Com relação ao item ‘ b ’, a maioria dos estudantes não chegou a responder.
Um número não desprezível de estudantes confundiu áreas funcionais da empresa com níveis hierárquicos, ou confundiu indicadores com ferramentas de gestão, como Canvas, Análise SWOT ou PDCA.
O desempenho geral ficou aquém do esperado.
A nota zero foi atribuída a 29 \% do total dos participantes.
As \textbf{repostas} em branco e as notas zero somaram quase 55 \% do total de respostas.
Dividindo as respostas apresentadas em três grupos, o resultado foi : respostas fracas ( notas entre zero e 30 ), corresponderam a 56 \% dos participantes ; respostas medianas ( notas entre 35 e 70 ) a cerca de 15 \%, e respostas boas ( notas entre 75 e 100 ), a apenas 3 \% do total.
O percentual de notas 100 chegou a 1 \% de todas as questões corrigidas.
As respostas consideradas boas trouxeram construções textuais mais bem cuidadas, com justificativas para a associação entre métodos tayloristas e a gestão da produção apresentadas de forma clara e correta.
Essas considerações eram complementadas pela apresentação de três indicadores de desempenho de três áreas distintas da gestão da produção.
As respostas medianas apresentavam, corretamente, a associação entre métodos tayloristas e a gestão da produção, mas não traziam explicações a respeito dos indicadores de desempenho. tayloristas e a gestão da produção, mas não traziam explicações a respeito dos indicadores de desempenho.
A correção parece revelar que há um certo desequilíbrio na formação dos administradores.
O fato de muitos estudantes terem respondido à questão de forma parcial, tratando apenas do primeiro item, pode significar que certos conteúdos não são priorizados na formação do administrador, ou que há dificuldades de aproximação entre conceitos teóricos e a prática.
O estudante parece ter mais facilidade em tratar de conteúdos puramente teóricos, como características da Administração Científica.
A dificuldade em propor indicadores de desempenho mostra essa desconexão. 6.3.2.3 - Análise da Questão Discursiva 4 do Componente de Conhecimento Específico Na Tabela 6.19, constam as informações relativas à questão 4 do conjunto de questões do Componente de Conhecimento Específico.
O desempenho dos estudantes nessa questão foi inferior ao das questões 3 e 5.
A Média geral do Brasil foi 14,2, sendo a maior Média registrada na região Sudeste ( 15,5 ), e a menor, na região Centro-Oeste ( 10,7 ).
A nota Máxima ( 100,0 ) foi atingida por, pelo menos, um concluinte de cada região.
A Mediana em todo o Brasil foi zero, o mesmo valor foi obtido nas cinco Grandes Regiões, indicando que, pelo menos, 50 \% dos participantes tiraram nota zero em cada uma das regiões.
Também foi zero a nota Mínima de todas as regiões.
No Gráfico 6.12, está representada a distribuição de notas da questão discursiva 4, do Componente de Conhecimento Específico.
Essa distribuição tem moda na classe de estudantes que, tendo respondido à questão, obtiveram nota 0,0, seguida da classe de alunos que deixaram a questão em branco.
A soma das frequências destes dois intervalos supera 63 \% dos respondentes.
O intervalo ( 10 ; 20 ] caracteriza -se como um máximo local, com cerca de 14 \% do total. 6.3.2.4 Comentários sobre as respostas à Questão Discursiva 4 A questão trata de tema relevante para a formação do administrador, especialmente marketing e relações com o mercado, contemplado pelas Diretrizes Curriculares do curso de Administração.
O texto introdutório, fragmento adaptado de uma obra reconhecida e atual, apresenta reflexões que auxiliavam o estudante a compreender aspectos relacionados com o tema e explicava o termo greenwashing, que não é propriamente corrente nos cursos superiores de Administração.
Trazia, também, alguns conceitos com relação ao pensamento de marketing e definia marketing sustentável que, ao mesmo tempo, está focado nas necessidades individuais dos clientes e nos compromissos das organizações com a responsabilidade social corporativa.
Ao estudante foi solicitado, no item ‘ a ’, explicar como os conceitos de marketing orientado para o cliente e marketing societal se complementam no conceito de marketing orientado ao interesse social.
O marketing orientado ao interesse social não é tratado comumente na literatura utilizada em cursos de graduação em administração, mesmo quando tratado como marketing sustentável ( outra nomenclatura encontrada no texto da questão ).
Essa dupla nomenclatura pode parecer um pouco confusa, mas o texto introdutório permite que o estudante reflita a respeito da orientação do marketing e apresente uma resposta ao item ‘ a ’.
No tocante ao item ‘ b ’, que solicitava que o estudante apresentasse e explicasse três exemplos de ações de greenwashing, as dificuldades foram maiores, o que talvez justifique o fato de muitos estudantes se dedicarem a responder apenas ao item ‘ a ’.
O padrão de resposta era bastante flexível e corretamente calibrado para avaliar o conhecimento dos estudantes nessa questão.
A princípio, considerou -se que essa questão tinha grau de dificuldade médio, em comparação às duas outras questões discursivas de Conhecimentos Específicos, por exigir que o estudante relacionasse teoria de administração, mais especificamente de marketing, e a prática das empresas.
Uma boa resposta exigia, ainda, atualização e acompanhamento da literatura e da mídia de negócios para poder \textbf{ultrapassar}, o que é mais comumente tratado na disciplina de marketing.
No entanto, essa foi a questão com o maior número de respostas em branco ( 28 \% ) e notas zero ( 35 \% ).
Assim, 63 \% dos estudantes ficaram com nota zero, ou por deixarem a questão em branco ou por não acertarem nenhum item solicitado, o que evidencia que a questão foi considerada difícil pelos estudantes.
A questão permitia que os estudantes refletissem a respeito de conceitos bastante atuais e manifestassem capacidade de associar corretamente teoria e prática, com liberdade para escolher três ações de greewashing no marketing.
Contudo, muitos apresentaram apenas repetições de trechos do enunciado em suas respostas e tiveram dificuldades em relacionar, corretamente, o marketing orientado ao cliente e o marketing societal.
Houve casos em que os estudantes relacionaram marketing social às redes sociais e trataram da expansão das redes sociais e suas supostas influências sobre os consumidores.
A necessidade de tratar do conceito de greenwashing e relacioná -lo a práticas empresariais se mostrou difícil para a maior parte dos estudantes.
O termo parece pouco conhecido e sua associação com a realidade do mercado e das práticas empresariais foi mencionada em poucas respostas.
No entanto, aqueles que conheciam o termo greenwashing trataram de casos como o recente escândalo de adulteração de testes da Volkswagen, ou de questões associadas à mineradora Samarco e à Vale, em Minas Gerais.
Por outro lado, muitos estudantes trataram greenwashing como sinônimo de qualquer prática desonesta ou antiética da empresa.
Por conta disso, foram listadas, como exemplos de greenwashin, práticas de preços que enganam o consumidor, falsas promoções, fraudes em garantias de produtos, atitudes contra a diversidade no ambiente de trabalho e exploração da mão de obra.
Houve aqueles que consideraram greenwashing uma prática positiva e citaram exemplos como uso de energia limpa, banimento de canudos plásticos ou sacolas plásticas.
O resultado dessa questão foi o pior entre as questões discursivas, se forem considerados o percentual de notas zero, 35 \% dos que tentaram responder à questão, e os que sequer responderam, 28 \%.
No entanto, essa foi também a questão que obteve o maior percentual de notas 100, cerca de 1,5 \%.
As respostas fracas, com notas entre zero e 30, corresponderam a 53 \% do total de participantes ; as respostas medianas ( notas entre 35 e 70 ) corresponderam a 14 \% ; e as boas ( notas entre 75 e 100 ) corresponderam a 4 \% do total de estudantes.
As \textbf{repostas} consideradas fracas trouxeram simples repetição do enunciado e/ou trataram, de forma incorreta, o marketing orientado para o social.
Também não trataram do greewashing e/ou realizaram associações equivocadas de práticas ao conceito de greenwashing no marketing.
As respostas medianas, em maior medida, foram dos estudantes que responderam corretamente ao item ‘ a ’, mas não apresentaram e explicaram três exemplos de greenwashing no marketing.
As respostas boas trouxeram explicações corretas para o marketing orientado para o social, assim como trouxeram e explicaram dois ou mais exemplos corretos de greewashing no marketing.
Assim como visto nas \textbf{repostas} dadas à Questão 3, as apresentadas a essa questão deixaram a impressão de haver dificuldade para estabelecer relação entre conceitos e práticas.
O enunciado trazia, de alguma forma, uma explicação a respeito do marketing orientado ao social, e o estudante poderia ter relacionado essa explicação, mesmo que breve, às práticas empresariais que ele observa no mercado.
No entanto, foram poucos os casos em que isso ocorreu.
Pode -se considerar que o uso de um conceito menos frequente na literatura mais tradicional de marketing, como greenwashing, pode ter dificultado mais ainda a ligação entre teoria e prática solicitada no item ‘ b ’ da questão.
De qualquer forma, ficou também a impressão de que os estudantes de administração não têm acompanhado a mídia de negócios, na qual o termo greenwashing é mais frequentemente encontrado.
Nesse ponto, aponta -se a possível necessidade de atualização do material didático utilizado para o ensino de marketing. 6.3.2.5 Análise da Questão Discursiva 5 do Componente de Conhecimento Específico Na Tabela 6.20, constam as informações relativas à questão 5 do conjunto do Componente de Conhecimento Específico.
O desempenho dos estudantes de todo o Brasil nesta questão foi o mais alto entre as três questões discursivas desse componente.
A nota Média dos estudantes de todo o Brasil foi 27,4.
A maior Média foi registrada na região Sudeste ( 28,5 ), enquanto a menor Média foi registrada na região Centro-Oeste ( 24,8 ).
Quanto Quanto à variabilidade das notas, o Desvio padrão dos alunos do Brasil, como um todo, foi 27,4.
Enquanto o maior desvio foi encontrado na região Sudeste ( 27,7 ), o menor foi encontrado na região Centro-Oeste ( 26,8 ).
A Mediana para o Brasil e em todas as regiões foi 25,0.
Para o conjunto de alunos de Administração do Brasil, a nota Máxima foi 100,0, obtida por, pelo menos, um aluno de todas as regiões.
A nota Mínima do Brasil foi zero, o mesmo valor encontrado em todas as regiões.
No Gráfico 6.13, é representada a distribuição das notas da questão discursiva 5 do Componente de Conhecimento Específico.
Essa distribuição tem modal o intervalo ( 40,50 ], correspondendo a cerca de 33 \% dos respondentes.
Aproximadamente 27 \% dos participantes tiraram zero nessa questão.
Somando -se a frequência desse intervalo com a daqueles que deixaram a questão em branco, chega -se próximo a 45 \% dos participantes. 6.3.2.6 Comentários sobre as respostas à Questão Discursiva 5 O enunciado da Questão 5 trazia um caso ilustrativo de vazamento de informações de clientes, muito relevante para os administradores.
As Diretrizes Curriculares Nacionais do Curso de Administração incluem os sistemas de informação e, portanto, a questão está alinhada ao conjunto de conhecimentos esperados de um formando da Área.
Em cursos de Administração, o tema de segurança de informação faz parte da ementa da disciplina Sistemas de Informação Gerencial ( SIG ).
Cursos de SIG comumente tratam das questões de segurança da informação com base na vulnerabilidade dos sistemas de informação, nas tecnologias, nas ferramentas para garantia da segurança das informações e nas formas de recuperação de desastres e planos de continuidade de negócios.
No entanto, questões conceituais a respeito de princípios de segurança da informação são raramente encontradas na literatura de SIG e são, portanto, pouco tratadas em cursos de Administração.
Dada a forma com que os negócios vêm se desenvolvendo na contemporaneidade, é essencial que esse tipo de conhecimento seja discutido por administradores e faça parte da formação desse profissional.
Considerou -se que essa questão pode ser qualificada como difícil porque, além de exigir que o estudante relacionasse conhecimento da administração a uma situação real ( relação entre teoria e prática ), abordava um conhecimento que é mais específico ( princípios de segurança da informação ).
No entanto, essa foi a questão com o menor índice de respostas em branco ( 18 \% ) e um percentual de notas zero comparável ao da Questão 3 ( 27 \% ), totalizando o menor índice de notas zero e o de respostas em branco somados ( 45 \% ).
O contexto do vazamento de dados é útil para reflexões que fazem parte do universo do administrador, no tocante à gestão de informações, ao marketing e ao comportamento ético das empresas.
O padrão de resposta era flexível e com foco na segurança da informação em termos mais próximos do que é tratado na disciplina de SIG.
Incluía, ainda, possíveis relações com posicionamento de mercado e com ética das organizações.
Por conta disso, o padrão conseguiu abranger diversas possibilidades de respostas, tratando o conhecimento do administrador como algo mais complexo do que a memorização de teorias ou ferramentas, sem ligação com a prática.
O comando da questão, “ redija um texto ”, apontou para a necessidade de o estudante construir um texto dissertativo, e isso pareceu problemático.
Muitos demonstraram dificuldades em escrever um texto dissertativo bem estruturado, que relacionasse o caso em tela com segurança da informação e posicionamento da empresa.
Grande parte dos textos das respostas apresentavam pouca fluidez, eram muito truncados, e algumas respostas foram apresentadas em itens ( com “ bullet points ” ).
O baixo conhecimento a respeito de segurança da informação fez com que muitos estudantes apenas retirassem trechos do enunciado e os reescrevessem em suas respostas.
A necessidade de redigir um texto, como solicitado no enunciado, fez com que a repetição de trechos do texto ilustrativo do enunciado servisse como estrutura do texto da resposta, mesmo em casos em que o estudante trazia alguma reflexão própria correta.
Isso ilustra a dificuldade de alguns estudantes em produzir um texto original bem estruturado, no qual consigam apresentar seus argumentos.
As respostas que receberam pontuação, na maioria das vezes, em algum grau, trataram de aspectos mercadológicos ou éticos associados ao vazamento de dados.
Aprofundamentos acerca da segurança de informação foram menos frequentes.
Foram poucos os estudantes que citaram corretamente os princípios de segurança da informação, e um número também pequeno conseguiu articular elementos da segurança da informação com a situação apresentada.
Essa questão apresentou, ao mesmo tempo, o menor percentual de notas zero e o menor percentual de notas acima de 80, entre as três questões discursivas : 27 \% e 0,3 \%, respectivamente.
Houve maior concentração de notas entre 40 e 75, em torno de 43 \%.
As respostas fracas ( diferentes de zero e com nota menor do que 35 ) corresponderam a 12 \%, e 17,7 \% sequer responderam à questão.
As \textbf{repostas} consideradas fracas se caracterizaram pela repetição de trechos do enunciado, sem nenhuma contribuição advinda de alguma reflexão do estudante.
Dessa forma, essas respostas não trouxeram nada ou muito pouco do que poderia ser considerado conhecimento a respeito das consequências práticas de problemas de segurança da informação.
A maioria desse grupo de respostas não trouxe menção alguma à segurança da informação, que não estivesse presente no enunciado.
Houve aquelas respostas genéricas, como “ a internet é um ambiente de risco ” ou “ o mundo digital é muito arriscado ”, que também foram consideradas fracas, além das respostas que tratavam de obviedades a respeito da responsabilidade da empresa e possíveis danos aos clientes, mas deixavam de lado conhecimentos pertinentes de algum campo da Administração.
As respostas medianas foram de dois tipos.
O primeiro tipo tratava da abordagem do posicionamento da empresa do ponto de vista mercadológico e oferecia alguma análise do caso, a partir desse ponto de vista.
O segundo tipo tratava de questões éticas relacionadas com vazamento de informações.
As respostas boas trataram de possíveis ações da empresa para reverter problemas de posicionamento e imagem, formas de recuperação de serviços utilizadas pela empresa e/ou de questões éticas relacionadas ao vazamento de informações, mas trataram, também, dos princípios de segurança da informação.
No geral, as respostas indicaram que os estudantes têm pouco conhecimento a respeito de segurança da informação.
Mesmo considerando apenas os conteúdos tradicionalmente ministrados na disciplina de Sistemas de Informação Gerencial, poucos estudantes construíram respostas razoáveis.
Houve dificuldade em articular a situação apresentada e formas de ação da empresa, mesmo com a possibilidade de o estudante tratar de questões mercadológicas ou éticas associadas ao vazamento de informações pessoais dos clientes.
Do ponto de vista estrutural, a questão solicitava que o estudante redigisse um texto com a análise do posicionamento da empresa citada e ficou clara a dificuldade com a produção textual.
Os textos apresentados eram, em sua maioria, pouco fluidos, truncados e de entendimento difícil. 6.3.3 Considerações Finais As questões discursivas de formação específica estavam alinhadas às Diretrizes Curriculares do Curso de Administração.
Tratavam de disciplinas importantes para a formação do administrador, sem deixar de lado a prática.
Na verdade, todas as questões desse exame buscavam fazer com que os estudantes precisassem relacionar o conhecimento conceitual com a realidade das empresas e do mercado e, portanto, avaliavam de forma objetiva a capacidade de atuação no mercado de trabalho.
A Questão discursiva 3 relacionava o conhecimento teórico mais tradicional da administração, a Administração Científica, com a Gestão da Produção e Operações e solicitava ao estudante que ele demonstrasse ser capaz de realizar a ligação entre áreas funcionais da empresa e indicadores de desempenho.
Esse conhecimento e essa habilidade são fundamentais para a atuação do administrador e, portanto, importantes para avaliar a formação desse profissional.
A Questão 4 pedia que o estudante fizesse uma reflexão mais avançada a respeito da orientação do marketing e apontasse para a necessidade de o administrador entender as relações entre mercado e sociedade.
Além disso, pedia que o estudante identificasse e explicasse ações de greenwashing no mercado atual.
A aproximação entre conhecimentos técnicos, como as ferramentas para atuação no mercado, e a responsabilidade social e ambiental são fundamentais para os administradores hoje em dia.
Dessa forma, a questão abordava, de maneira direta, um ponto importante para o desenvolvimento de profissionais mais conscientes de suas ações dentro das empresas.
A Questão 5 trazia um tema tratado de maneira ainda tímida na administração, que é a segurança de informações.
Avaliava a capacidade de o estudante encontrar subsídios teóricos para sua resposta ( princípios da segurança da informação ) e tratar do posicionamento da empresa frente a um caso ilustrativo muito interessante e atual de vazamento de dados de clientes.
O tema da tecnologia da informação se mostra cada vez mais relevante para os administradores, e essa questão trazia esse tema de forma muito pertinente.
As respostas dos estudantes, de forma geral, demonstraram a dificuldade em relacionarem teoria e prática, competência exigida em todas as questões.
Além disso, ficou aparente o reduzido conhecimento dos estudantes a respeito do que está ao seu redor em termos de informações sobre negócios, mercados e gestão de empresas na prática.
Mais especificamente, no tocante ao conteúdo, o conhecimento a respeito de Tecnologia da Informação, tão em voga nos dias de hoje, mostra -se limitado entre os estudantes de administração.
Por último, as respostas apresentadas pelos estudantes também mostraram a dificuldade que a maioria dos estudantes tem de produzir textos de forma articulada e bem estruturada, o que denota que, mesmo sendo deficiências de formação prévia, os cursos de administração talvez não as estejam conseguindo amenizá -las.

% matched lemmas: aprendizado, centrado, chão, destruir, esporádico, mentira, observação, personificar, peso, repor, sócio-histórico-cultural, ultrapassar, étnico-racial
\end{textsample}
